\section{Two large path counts in a three-terminal graph}
\label{sec:three-pole-high-pair-tradeoff}

A \emph{normalized three-pole} is a connected loopless multigraph with three distinct degree-two terminals.
Every other vertex has degree three.
Its rank $b$ and order $n$ satisfy $n=2b+1$.
Paths distinguish parallel edges.

\begin{theorem}[Two high pairs force a small third pair]
\label{thm:three-pole-two-high}
For every $0<\varepsilon\le1/2$, there is a constant $Q_\varepsilon$ such that every normalized three-pole $T$ satisfies
\[
 P_{ac}(T),P_{bc}(T)\ge(1/2+\varepsilon)2^b
 \quad\Longrightarrow\quad
 P_{ab}(T)\le Q_\varepsilon 2^{b/2}.
\]
One may take $Q_\varepsilon=2^{30C_\varepsilon+20}$, where $C_\varepsilon$ is the additive path-stability constant.
For a simple three-pole, $2^{18C_\varepsilon+20}$ suffices.
\end{theorem}

\begin{proof}
We first explain the use of additive stability in a normalized multigraph.
Suppress the third terminal $b$ when considering the pair $a,c$.
Its two neighbors are distinct.
Otherwise, its pendant digon would contribute rank while avoiding every $a$--$c$ path.
This would bound the normalized path count by one half.
The resulting graph is a normalized two-pole $N$ of the same rank and path count.

The additive stability proof applies directly to every shortest terminal path in $N$.
Indeed, subdividing parallel edges gives a simple graph whose normalized reduction returns $N$.
No cubic expansion is necessary.
Thus its terminal distance is at least $b-C_\varepsilon$.
The same lower bound holds in $T$.
Repeating the argument for the pair $b,c$ gives
\[
 d:=\min\{\operatorname{dist}_T(c,a),\operatorname{dist}_T(c,b)\}
 \ge b-C_\varepsilon.
\]

We will also need a shortest $c$--$a$ path $A$ in $T$ with
\[
 \beta(T-V(A))\le C_\varepsilon.
\]
Let $e$ be the suppressed edge in $N$.
If $\operatorname{dist}_T(c,a)=\operatorname{dist}_N(c,a)$, choose a shortest path in $N$ that avoids $e$.
Such a path exists, since any path using $e$ gains one edge when lifted.
Otherwise, the two distances differ by one, and every shortest path in $N$ lifts to a shortest path in $T$.
In either case, the residual-rank conclusion of additive stability applies to the chosen path in $N$.
Its lift has the same residual rank, since the extra vertex only subdivides an edge or becomes a leaf or isolated vertex.
This proves the required bound.

Assume first that $T$ is bridgeless.
It has no cut vertex, since it is subcubic.
Let $L_i$ be its distance layers from $c$.
Every layer $L_i$ with $1\le i<d$ separates $c$ from both other terminals, so $|L_i|\ge2$.
The terminal tail $\bigcup_{i\ge d}L_i$ contains both $a,b$.
Since $n=2b+1$, it follows that $d\le b$.
Put $\Delta=b-d\le C_\varepsilon$.

Mark $c$, the whole terminal tail, and every internal layer of size at least three.
The baseline counts are one vertex at $c$, two in each internal layer, and two in the terminal tail.
Their sum is $2d+1$, leaving excess $2\Delta$.
Using $w\le3(w-2)$ for $w\ge3$, the marked set has size at most $6\Delta+3$.
Every remaining layer has exactly two vertices.
It contains one vertex of $A$ and one vertex outside $A$.

In the multigraph case, mark every whole layer containing an endpoint of a parallel pair outside $A$.
There are at most $H:=\beta(T-V(A))$ such pairs, since their digons are independent.
These extra markings add at most $4H$ vertices.
A parallel pair cannot have just one endpoint on $A$, since each path vertex has at most one incidence outside $A$.
A parallel pair with both endpoints on $A$ joins consecutive path vertices, by geodesicity.

Fix a simple $a$--$b$ path $P$.
If $P$ stays in the terminal tail, its entire edge set lies in the marked set.
Otherwise, let $i<d$ be its minimum distance layer.
If that layer is not already marked, mark its two vertices.
The final marked set $B$ therefore satisfies
\[
 |B|\le6\Delta+4H+5.
\]
Fix all edges of $P$ incident with $B$, with at most $2^{3|B|}$ choices.

Every unmarked layer strictly between $i$ and $d$ carries two strands of $P$.
Indeed, both path endpoints lie beyond that layer, so every prefix cut meets $P$ evenly.
Connectedness forces both layer vertices to carry one backward and one forward path edge.
A same-layer path edge would give an odd prefix crossing.
Separate backward and forward caps would disconnect the path.

Across consecutive unmarked layers, these strands therefore form a perfect matching.
The matching is unique unless the cut is $K_{2,2}$ or contains a doubled edge of $A$.
Each exceptional cut has at most two choices after the extra markings.
Two such cuts cannot be consecutive, since their shared path vertex would have degree at least four.
Thus their combined contribution is at most $2^{(d-i)/2}$.
All remaining edges are fixed by the marked-edge data.

Summing over the minimum layer, and including paths wholly in the tail, gives
\[
 \begin{split}
 P_{ab}(T)
 &\le 2^{18\Delta+12H+15}
       \left(1+\sum_{i=0}^{d-1}2^{(d-i)/2}\right)\\
 &< 5\cdot2^{18\Delta+12H+15+d/2}
 \le 2^{30C_\varepsilon+20+b/2}.
 \end{split}
\]
For a simple graph, no parallel-edge markings are needed, so the term $12H$ is omitted.

Finally suppose that $T$ has a bridge.
Every bridge shore contains a cycle.
Indeed, a shore with $m$ vertices and $p$ original terminals has rank $(m-p+1)/2$, a positive integer.
The two high path counts force every bridge to separate $\{a,b\}$ from $c$.
No original terminal can be a bridge endpoint, since the corresponding terminal path would immediately leave a cyclic shore.

Delete one bridge, and retain the shore containing $a,b$.
It is a normalized three-pole with third terminal at the bridge attachment.
Its two normalized path densities remain at least $1/2+\varepsilon$, by factorization across the bridge.
Every $a$--$b$ path stays in this shore.
Its rank is smaller than $b$.
Induction on rank proves the same estimate, with the same constant.
\end{proof}

\begin{corollary}[A large-shore three-edge cut]
\label{cor:large-three-cut-half-exponent}
There are absolute constants $B_0,K$ such that the following holds.
Let $G$ be a connected cubic graph with a three-edge cut whose two shores are normalized three-poles of ranks at least $B_0$.
Then, writing $r=\beta(G)$,
\[
 c(G)\le2^{r-1}+K2^{r/2}.
\]
Parallel edges are allowed.
\end{corollary}

\begin{proof}
Write the shores as $T_1,T_2$, with ranks $b_1,b_2$.
Index each shore's terminal pairs by the same three pairs of cut edges.
Put
\[
 f_{ij}=\frac{P_j(T_i)}{2^{b_i}},\qquad
 f_i=(f_{i1},f_{i2},f_{i3}),\qquad c_i=c(T_i).
\]
Each cycle crosses the cut either zero or two times, so
\[
 r=b_1+b_2+2,\qquad
 c(G)=c_1+c_2+2^{b_1+b_2}\langle f_1,f_2\rangle.
\]

Call a coordinate high when it is at least $2/3$.
Put $Q=Q_{1/6}$ and choose
\[
 B_0\ge\max\{5,\lceil2\log_2(16Q)\rceil\}.
\]
The theorem shows that a shore with two high coordinates has third coordinate at most
$\delta_i:=Q2^{-b_i/2}\le1/16$.
Consequently, no shore has three high coordinates.

If both shores have two high coordinates in the same positions, then
\[
 \langle f_1,f_2\rangle\le2+\delta_1\delta_2.
\]
In every other case the dot product is at most $17/9$.
For two shores with one high coordinate each, this follows from $1+2(2/3)^2=17/9$.
If exactly one shore has two high coordinates, the bound is at most $1+2/3+(2/3)/16<17/9$.
If both have two high coordinates in different positions, the bound is at most $1+2/16<17/9$.
A shore with no high coordinate gives a smaller bound in every case.

If each shore has at least one high coordinate, suppress its unused terminal.
Theorem~\ref{thm:high-path-cycle-tradeoff} then gives
\[
 c_i\le K_{1/6}2^{b_i/2}.
\]
The aligned case has crossing excess at most $Q^2 2^{(b_1+b_2)/2}$.
The other cases have no crossing excess.
Since $b_1,b_2\ge2$, their internal cycle counts add to at most
$K_{1/6}2^{(b_1+b_2)/2}$.
This gives the claimed estimate.

If one shore has no high coordinate, the crossing deficit is at least
$2^{b_1+b_2}/9$.
The trivial internal bounds satisfy
\[
 c_1+c_2<2^{b_1}+2^{b_2}
 \le2^{b_1+b_2}/16,
\]
because $b_1,b_2\ge5$.
Thus $c(G)<2^{r-1}$ in this case.
Taking $K=Q^2+K_{1/6}$ proves the corollary.
\end{proof}

\begin{theorem}[The three-terminal path sum]
\label{thm:universal-tripole-sum}
There is an absolute constant $S$ such that every normalized three-pole $T$ of rank $b$ satisfies
\[
 P_{ab}(T)+P_{ac}(T)+P_{bc}(T)
 \le2^{b+1}+S2^{b/2}.
\]
One may take $S=\max\{6,Q_{1/32}\}$.
\end{theorem}

\begin{proof}
If one terminal pair is full, Lemma~\ref{lem:one-full-tripole-sum} proves the estimate with constant six.
Otherwise, Theorem~\ref{thm:terminal-path-saturation-gap} bounds every count by $(15/16)2^b$.
If their sum is at most $2^{b+1}$, the claim is immediate.
If their sum exceeds $2^{b+1}$, at least two counts exceed $(17/32)2^b$.
Indeed, one count at most $15/16$ and two counts at most $17/32$ have normalized sum at most two.
Apply Theorem~\ref{thm:three-pole-two-high} with $\varepsilon=1/32$.
It bounds the remaining count by $Q_{1/32}2^{b/2}$.
The other two counts are each at most $2^b$, proving the theorem.
\end{proof}

\begin{theorem}[Every cyclic cut of size at most three]
\label{thm:all-three-cut-half-exponent}
There is an absolute constant $K_3$ such that every connected loopless cubic multigraph $G$ with a cyclic edge cut of size at most three satisfies
\[
 c(G)\le2^{r-1}+K_3 2^{r/2},\qquad r=\beta(G).
\]
Here a cut is cyclic when both shores contain a cycle.
In a cubic graph, every nontrivial three-edge cut is cyclic.
\end{theorem}

\begin{proof}
A graph with a bridge has $c(G)\le2^{r-1}$.
A two-edge cut is covered by Corollary~\ref{cor:cubic-multigraph-two-cut-error}.
It remains to consider a three-edge cut in a graph with no smaller edge cut.

Each shore is connected.
Otherwise, its separate components would give two cuts of size at least three within the specified three-edge boundary.
The three attachment vertices on each shore are distinct.
If one vertex had two boundary edges, deleting that vertex from its shore would expose a two-edge cut.
Three boundary edges at one vertex would make that shore a singleton, contrary to cyclicity.
Thus the shores $T_1,T_2$ are normalized three-poles.
Their ranks $b_1,b_2$ are positive, and $r=b_1+b_2+2$.

Let $B_0$ be the constant from Corollary~\ref{cor:large-three-cut-half-exponent}.
If both ranks are at least $B_0$, that corollary applies.
Otherwise, reorder the shores so that $1\le b_1<B_0$ and $b_1\le b_2$.

Let $R_0$ be an integer such that every graph of rank at least $R_0$ has at most $(9/16)2^b$ cycles.
The established asymptotic cycle bound supplies such an integer.
Subdividing parallel edges extends that bound to loopless multigraphs.
Put
\[
 M=\max\{R_0,B_0+3\}.
\]
If $b_2<M$, then $r\le2M$.
The trivial cycle bound is absorbed by the constant $2^{M+1}$ in the asserted estimate.
Assume henceforth that $b_2\ge M$.

Suppose first that one terminal pair of $T_2$ has at least $(9/16)2^{b_2}$ paths.
Suppressing its unused terminal and applying Theorem~\ref{thm:high-path-cycle-tradeoff} gives
\[
 c(T_2)\le K_{1/16}2^{b_2/2}.
\]
Every path count in $T_1$ is at most $2^{b_1}$.
Theorem~\ref{thm:universal-tripole-sum} therefore bounds the crossing cycles by
\[
 2^{b_1}\sum_jP_j(T_2)
 \le2^{r-1}+S2^{b_1+b_2/2}.
\]
Since $b_1<B_0$, the last term is at most
$S2^{B_0/2}2^{(b_1+b_2)/2}$.
Also $c(T_1)<2^{b_1}\le2^{(b_1+b_2)/2}$.
Together, these estimates give the required bound with constant
$S2^{B_0/2}+K_{1/16}+1$.

Finally suppose that every terminal-pair count in $T_2$ is below $(9/16)2^{b_2}$.
The crossing cycles then number at most
\[
 \frac{27}{16}\,2^{b_1+b_2}.
\]
Their deficit from $2^{r-1}$ is at least
\[
 \frac5{16}2^{b_1+b_2}\ge\frac58\,2^{b_2}.
\]
The choice of $R_0$ gives $c(T_2)\le(9/16)2^{b_2}$.
The remaining deficit is at least $2^{b_2}/16$.
Since $b_2\ge B_0+3\ge b_1+4$, it absorbs $c(T_1)<2^{b_1}$.
Thus $c(G)<2^{r-1}$ in this case.

For explicit dependence, let $K_2$ be the two-edge-cut constant and let $K_{\mathrm{large}}$ be the large-shore constant.
The choice
\[
 K_3=\max\left\{
 K_2,K_{\mathrm{large}},2^{M+1},
 S2^{B_0/2}+K_{1/16}+1
 \right\}
\]
suffices.
\end{proof}

This result leaves cyclically four-edge-connected cubic kernels as the remaining unrestricted case.
It does not prove the half-exponent bound for that case.
No publication-priority claim is made here.
